\documentclass[10pt,a4paper]{amsart}
\usepackage[margin=19mm]{geometry}
\usepackage{amsmath,amssymb,mathtools}
\usepackage[scr=rsfs,cal=euler]{mathalfa}
\usepackage{xfrac}
\usepackage[unicode,hidelinks,bookmarks=false]{hyperref}
\newcommand{\ie}{i.e.\ }
\newcommand{\cf}{cf.\ }
\newcommand{\resp}{respectively\ }
\newcommand{\Subsection}{\subsection}
\providecommand{\nobreakdash}{\nobreak\hbox{-}}
\setlength{\emergencystretch}{2em}
\multlinegap0pt
\usepackage{bm}
\usepackage{bbm}
\usepackage{dsfont}
\usepackage{xspace}
\usepackage{cancel}
\usepackage{mathtools}
\usepackage{amsthm}
\makeatletter
\@ifundefined{theorem}{\newtheorem{theorem}{Theorem}}{}
\@ifundefined{lemma}{\newtheorem{lemma}[theorem]{Lemma}}{}
\makeatother
\newcommand{\J}{\mathcal{J}}
\newcommand{\Kl}{\mathcal{K}}
\newcommand{\abs}[1]{\left|{#1}\right|}
\newcommand{\sq}{\mathfrak{s}}
\newcommand{\val}{{\operatorname{val}}}
\renewcommand{\O}{\mathcal{O}}
\title[On the Cubic Shimura lift to $PGL(3)$: The Fundamental Lemma]{On the Cubic Shimura lift to $PGL(3)$: The Fundamental Lemma}
\author{Solomon Friedberg, Omer Offen}
\date{Source: arXiv:2202.01247v2 (CC BY 4.0)}
\begin{document}
\maketitle

\noindent\emph{Source and licence.} On the Cubic Shimura lift to $PGL(3)$: The Fundamental Lemma. Version arXiv:2202.01247v2 (CC BY 4.0), distributed under the Creative Commons Attribution 4.0 International licence, as recorded in the arXiv licence metadata for this exact version and shown on the abstract page https://arxiv.org/abs/2202.01247v2. Full rights notice: licenses/arxiv-2202.01247-CC-BY-4.0.txt.\par\smallskip

\noindent\textit{Editorial scope.} Counted unit: the Lemma whose source label is \texttt{lem j>0} in the article (source0.tex lines 2120--2130, source label \texttt{lem j>0}), delivered together with the article's own complete original proof of it (source0.tex lines 2131--2200, one \texttt{proof} environment of 70 lines). No mathematics is written, repaired or re-derived: statement and proof are the article's own text line for line. Required context retained uncounted: lem kl int (lines 1257--1266); eq j>0 (lines 2042--2044). Every definition, display and label of those lines is retained unchanged. Omitted from this package and not redistributed: 1--1256, 1267--2041, 2045--2119, 2201--3417. Nothing inside a retained excerpt is omitted or altered: every retained body is delivered exactly as the article's lines are written, so no omission receipt of any kind accompanies this package. Citations: the retained excerpts carry no citation command at all, so no bibliography entry of the article is transcribed and no reference is redistributed. Reference adaptation: none was needed and none was made; every reference inside the delivered unit resolves inside it, so no unresolved reference or citation is produced. Printed slips kept literal: no printed word, symbol, hypothesis or number of the counted statement or of its proof was altered. Layout and class: the article's own class is \texttt{amsart} with packages including amssymb, a4wide, mathdots, url, hyperref, graphicx, color, enumerate; the wrapper is the lane's pinned \texttt{amsart} editorial wrapper at 10pt on A4, which restates the theorem-like environments the retained text needs and adds the article's own macro definitions that the retained text needs (\texttt{\textbackslash J}, \texttt{\textbackslash Kl}, \texttt{\textbackslash O}, \texttt{\textbackslash abs}, \texttt{\textbackslash sq}, \texttt{\textbackslash val}), with extra wrapper packages (bm, bbm, dsfont, xspace, cancel, mathtools). The delivered numbering is the unit's own. Assets: no retained span contains a figure environment, a graphics-inclusion command, a PDF-page inclusion, a table or a TikZ picture; no image file is copied into this package, the package declares no asset dependency and no image file is redistributed with it. Licence: version arXiv:2202.01247v2 of this article is distributed under the Creative Commons Attribution 4.0 International licence, as recorded in the arXiv licence metadata for this exact version and shown on the abstract page https://arxiv.org/abs/2202.01247v2. A full rights notice is delivered at licenses/arxiv-2202.01247-CC-BY-4.0.txt. Characters: every retained excerpt is pure ASCII, so the delivered engine, XeLaTeX with its default Latin Modern text font, reports no missing character.
\begin{lemma}\label{lem kl int}
We have
\[
\Kl(y;a,b)=\begin{cases} (1-q^{-1})\delta_{3\mid \val(y)} & \max(\abs{a},\abs{b})\le 1 \\
q^{-1}\sum_{u\in k_F^*}(y,u)_3\,\psi(au+bu^{-1}) &  \max(\abs{a},\abs{b})=q \\
q^{-\lfloor\frac{1-\val(a)}2\rfloor}(y,ab^{-1})_3\,\sum_{\epsilon=\pm 1} \psi(2\epsilon\sqrt{ab})\sq(\epsilon\sqrt{ab})& \max(\abs{a},\abs{b})>q \text{ and }a^{-1}b\in \O^{*2} \\
0 & \max(\abs{a},\abs{b})>q \text{ and }a^{-1}b\not\in \O^{*2}.
 \end{cases}
\]
\end{lemma}

\begin{multline}\label{eq j>0}
\J_{>0}(a,b)=\abs{a^{-1}b}\int_{D'_{>0}[a,b]} (x_1,bx_1y_2-1)_3(az_2,b)_3\\ \psi(-(b z_2)^{-1}(1-bx_1y_2)^{-1}x_1+y_1+x_2+az_2y_2)\ dx_1\ dx_2\ dy_1\ dy_2\ dz_2
\end{multline}

\begin{lemma}\label{lem j>0}
Let $\abs{b}\le \abs{a}< 1$. Then
\[
\J_{>0}(a,b)=\sum_{\ell=\lfloor \frac{\val(a)+1}2\rfloor}^{\val(a)-1}\J_\ell(a,b)
\]
where
\begin{multline*}
\J_\ell(a,b)=\abs{a}^{-2}\int_{\abs{y_2}=q^{-2\ell}\abs{b}^{-1},\,\abs{1-y_2}=\abs{ab^{-1}}} (b^{-1}\varpi^\ell,y_2-1)_3\\ \Kl(b;-ab^{-2}\varpi^\ell(1-y_2)^{-1},\varpi^{-\ell} y_2)\Kl(b^{-1}(y_2-1)^{-1};a^{-1}\varpi^\ell ,a^{-1}by_2\varpi^{-\ell} )\ dy_2.
\end{multline*}
In particular, $\J_{>0}(a,b)=0$ if $\abs{a}=q^{-1}$.
\end{lemma}

\begin{proof}
Note that the domain $D_{>0}[a,b]$ is characterized by the conditions
\[
\abs{ab^{-1}}<\abs{x_1},\abs{y_2}<\abs{b}^{-1};\, az_2\in \O^*;\, y_1-a^{-1}by_2,\,x_2-a^{-1}bx_1\in \O;\, \abs{1-bx_1y_2}=\abs{ab^{-1}}.
\]
In particular, if $\abs{a}=q^{-1}$ then $D_{>0}[a,b]$ is empty and $\J_{>0}(a,b)=0$. For the rest of this section assume that $\abs{a}<q^{-1}$.
Applying the variable changes $y_1\mapsto y_1+a^{-1}by_2$ and $x_2\mapsto x_2+a^{-1}bx_1$ to \eqref{eq j>0},
the integrand becomes independent of $y_1,\,x_2\in\O$. After integrating over these two variables we obtain the expression
\begin{multline*}
\J_{>0}(a,b)=\abs{a^{-1}b}\int (x_1,bx_1y_2-1)_3  (az_2,b)_3\\
 \psi((a^{-1}b-(b z_2)^{-1}(1-bx_1y_2)^{-1})x_1+(a^{-1}b+az_2)y_2)\ dx_1\ dy_2\ dz_2
\end{multline*}
where integration is over the domain defined by the conditions
\[
az_2\in\O^*,\ \ \ \abs{ab^{-1}}<\abs{x_1},\abs{y_2}<\abs{b}^{-1}\ \ \  \text{and} \ \ \ \abs{1-bx_1y_2}=\abs{ab^{-1}}. 
\]

Next, applying the variable change 
\[
x_1\mapsto az_2x_1, \ y_2\mapsto (az_2)^{-1}y_2
\]
we have
\begin{multline*}
\J_{>0}(a,b)=\abs{a^{-1}b}\int (x_1,bx_1y_2-1)_3  (az_2,b(bx_1y_2-1))_3 \\ \psi(bz_2x_1-ab^{-1}(1-bx_1y_2)^{-1}x_1+a^{-2}by_2z_2^{-1}+y_2)\ dx_1\ dy_2\ dz_2
\end{multline*}
where integration is over the same domain.
Now applying the variable change 
\[
z_2\mapsto z_2x_1^{-1},\,y_2\mapsto y_2(bx_1)^{-1}
\]
we have
\begin{multline*}
\J_{>0}(a,b)=\abs{a}^{-1}\int \abs{x_1}^{-2}(b,x_1)_3  (az_2,b(y_2-1))_3 \\ \psi(bz_2-ab^{-1}(1-y_2)^{-1}x_1+a^{-2}y_2z_2^{-1}+b^{-1}y_2x_1^{-1})\ dx_1\ dy_2\ dz_2
\end{multline*}
where integration is over the domain defined by
\[
\abs{ab^{-1}}<\abs{x_1}=\abs{az_2}<\abs{b}^{-1},\ \ \ \abs{ax_1}<\abs{y_2}<\abs{x_1}\ \ \  \text{and} \ \ \ \abs{1-y_2}=\abs{ab^{-1}}. 
\]
Note further that the condition $\abs{1-y_2}=\abs{ab^{-1}}$ implies that $\abs{y_2}\le \abs{ab^{-1}}$ and combined with the condition $\abs{ab^{-1}}<\abs{x_1}$ it 
implies that $\abs{y_2}<\abs{x_1}$. Thus the condition $\abs{y_2}<\abs{x_1}$ may be omitted.

We express the integral $\J_{>0}(a,b)$ as a sum over $\ell=1,\dots,\val(a)-1$ of integrals over the subdomain 
where $\abs{x_1}=\abs{b}^{-1}q^{-\ell}$, that is, over the domain defined by the conditions
\[
\abs{x_1}=\abs{b}^{-1}q^{-\ell},\ \abs{z_2}=\abs{ab}^{-1}q^{-\ell},\ q^{-\ell}\abs{b^{-1}a}<\abs{y_2}\ \ \ \text{and}\ \ \ \abs{1-y_2}=\abs{ab^{-1}}.
\]
Applying the variable change
\[
x_1\mapsto b^{-1}\varpi^\ell x_1,\ z_2\mapsto a^{-1}b^{-1}\varpi^\ell z_2
\]
to the $\ell$-th summand we have
\begin{multline*}
\J_{>0}(a,b)=\abs{a}^{-2}\sum_{\ell=1}^{\val(a)-1}\int (b^{-1}\varpi^\ell,y_2-1)_3\int_{\O^*}\int_{\O^*}(b,x_1)_3  (z_2,b(y_2-1))_3  \\ \psi(a^{-1}\varpi^\ell z_2-ab^{-2}\varpi^\ell(1-y_2)^{-1} x_1+a^{-1}b \varpi^{-\ell} y_2 z_2^{-1}+\varpi^{-\ell} y_2x_1^{-1})\ dx_1\ dz_2 \ dy_2\\=
 \abs{a}^{-2}\sum_{\ell=1}^{\val(a)-1}\int  (b^{-1}\varpi^\ell,y_2-1)_3\\ \Kl(b;-ab^{-2}\varpi^\ell(1-y_2)^{-1},\varpi^{-\ell} y_2)\Kl(b^{-1}(y_2-1)^{-1};a^{-1}\varpi^\ell ,a^{-1}by_2\varpi^{-\ell} )\ dy_2
\end{multline*}
where in the $\ell$-th integral $y_2$ is integrated over the domain defined by the conditions
\[
q^{-\ell}\abs{b^{-1}a}<\abs{y_2}\ \ \ \text{and}\ \ \ \abs{1-y_2}=\abs{ab^{-1}}.
\]

It follows from Lemma \ref{lem kl int} that in the domain of integration above we have 
\[
\Kl(b^{-1}(y_2-1)^{-1};a^{-1}\varpi^\ell ,a^{-1}by_2\varpi^{-\ell} )=0
\] 
unless 
\[
\abs{y_2}=q^{-2\ell}\abs{b}^{-1}.
\] 
Indeed, if $\ell<\val(a)-1$ then $\abs{a^{-1}\varpi^\ell}>q$ and for $\ell=\val(a)-1$ we have $\abs{a^{-1}\varpi^\ell}=q$ while in the domain of integration $\abs{a^{-1}by_2\varpi^{-\ell}}\ge q$. Since in the domain of the $\ell$-th integral we have $\abs{y_2}\le \abs{ab^{-1}}$ and since $q^{-2\ell}\le \abs{a}$ if and only if $\ell\ge \lfloor \frac{\val(a)+1}2\rfloor$ the lemma follows.
\end{proof}

\end{document}
